% =============================================================================
% APPENDIX XII: DOMAIN-FOUND - Domain Foundations of Complementarity
% =============================================================================
% Category: DOMAIN (Application Domains)
% Code: XII
% Version: 1.0 (January 2026)
%
% Purpose: Bridge between pure mathematical theory (Appendix X) and
%          domain-specific applications. Establishes how abstract
%          complementarity manifests in Economics, Psychology, and Management.
%
% Dependencies:
%   - Appendix X (FORMAL-FOUND): Mathematical foundations
%   - Appendix XI (FORMAL-MODELS): Canonical mathematical models
%   - Appendix VII (FORMAL-GTC): General Theory of Complementarity
%
% Layer Position: Layer 3 in the 4-Layer Architecture
%   Layer 1: Pure Mathematics (X)
%   Layer 2: Mathematical Models (XI)
%   Layer 3: Domain Foundations (XII) ← THIS APPENDIX
%   Layer 4: EBF Application (XIII)
%
% =============================================================================

\chapter{Domain Foundations of Complementarity}
\label{app:domain-found}

\begin{abstract}
This appendix bridges the gap between pure mathematical complementarity theory
(Appendix X) and domain-specific applications. We establish how the abstract
mathematical structures—lattices, supermodularity, and interaction tensors—manifest
in three primary domains: \textbf{Economics} (utility theory, production, games),
\textbf{Psychology} (behavioral interactions, cognitive complementarities), and
\textbf{Management} (organizational design, strategic fit). For each domain, we
provide: (1) the interpretation of $\gamma$, (2) domain-specific axioms,
(3) measurement approaches, and (4) typical parameter ranges from empirical literature.
This creates the foundation for Layer 4: EBF-specific application.
\end{abstract}

% ============================================================================
% QUICK REFERENCE
% ============================================================================
\begin{tcolorbox}[colback=blue!5!white,colframe=blue!75!black,title=\textbf{Quick Reference: Layer 3 Position}]
\textbf{The 4-Layer Architecture:}
\begin{enumerate}[noitemsep]
    \item \textbf{Layer 1 (Appendix X):} Pure mathematics — domain-independent axioms, lattice theory
    \item \textbf{Layer 2 (Appendix XI):} Mathematical models — six canonical forms (A-F)
    \item \textbf{Layer 3 (This Appendix):} Domain foundations — Economics, Psychology, Management
    \item \textbf{Layer 4 (Appendix XIII):} EBF application — behavioral change, intervention design
\end{enumerate}

\textbf{Key Translation:}
\begin{center}
\begin{tabular}{lll}
\textbf{Math (X/XI)} & $\rightarrow$ & \textbf{Domain (XII)} \\
\hline
Set $S$ & $\rightarrow$ & Choice set, action space, strategy set \\
Function $f: 2^S \to \mathbb{R}$ & $\rightarrow$ & Utility, production, payoff function \\
$\gamma_{ab} = f(\{a,b\}) - f(\{a\}) - f(\{b\}) + f(\emptyset)$ & $\rightarrow$ & Synergy, complementarity, fit \\
Lattice $(L, \leq)$ & $\rightarrow$ & Preference order, technology order \\
\end{tabular}
\end{center}
\end{tcolorbox}

% ============================================================================
% PART I: THE TRANSLATION PROBLEM
% ============================================================================
\section{The Translation Problem: From Mathematics to Domains}
\label{sec:translation}

\subsection{Why Domain Translation Matters}

\begin{tcolorbox}[colback=yellow!5!white,colframe=yellow!75!black,title=\textbf{The Central Challenge}]
Pure mathematical complementarity (Appendix X) is \textit{domain-agnostic}:
\begin{equation}
\gamma_{ab} = f(\{a,b\}) - f(\{a\}) - f(\{b\}) + f(\emptyset)
\end{equation}
But application requires \textit{domain-specific} interpretation:
\begin{itemize}[noitemsep]
    \item What are $a$ and $b$? (goods, behaviors, practices, beliefs?)
    \item What is $f$? (utility, production, fitness, payoff?)
    \item What does $\gamma > 0$ mean? (synergy, fit, enhancement?)
    \item How is $\gamma$ measured? (experiments, surveys, revealed preference?)
\end{itemize}
\end{tcolorbox}

\subsection{The Translation Map}

\begin{definition}[Domain Translation]
A \textbf{domain translation} $\mathcal{T}_D$ maps the mathematical structure $(S, f, \gamma)$ to a domain-specific interpretation:
\begin{equation}
\mathcal{T}_D: (S, f, \gamma)_{math} \mapsto (X_D, U_D, \gamma_D)_{domain}
\end{equation}
where:
\begin{itemize}[noitemsep]
    \item $X_D$ is the domain-specific element space
    \item $U_D$ is the domain-specific value function
    \item $\gamma_D$ is the domain-specific complementarity interpretation
\end{itemize}
\end{definition}

\subsection{Three Primary Domains}

\begin{table}[h]
\centering
\caption{Domain Translation Overview}
\label{tab:domain-overview}
\renewcommand{\arraystretch}{1.4}
\begin{tabular}{@{}p{2.5cm}p{3cm}p{3cm}p{4cm}@{}}
\toprule
\textbf{Domain} & \textbf{Elements $X_D$} & \textbf{Value $U_D$} & \textbf{$\gamma_D$ Interpretation} \\
\midrule
\textbf{Economics} & Goods, actions, strategies & Utility, profit, payoff & Consumption synergy, strategic complement \\
\textbf{Psychology} & Behaviors, cognitions, traits & Well-being, performance & Behavioral interaction, cognitive fit \\
\textbf{Management} & Practices, resources, capabilities & Firm value, performance & Organizational fit, strategic alignment \\
\bottomrule
\end{tabular}
\end{table}

% ============================================================================
% PART II: ECONOMIC FOUNDATIONS
% ============================================================================
\section{Economic Foundations of Complementarity}
\label{sec:economics}

\subsection{Historical Development}

\begin{tcolorbox}[colback=gray!5!white,colframe=gray!75!black,title=\textbf{Timeline: Complementarity in Economics}]
\begin{description}[noitemsep]
    \item[1838] Cournot: First analysis of complementary goods
    \item[1881] Edgeworth: Complementarity in utility functions
    \item[1934] Hicks \& Allen: Cross-price elasticity definition
    \item[1947] Samuelson: Revealed preference foundations
    \item[1978] Topkis: Supermodularity in optimization
    \item[1990] Milgrom \& Roberts: Strategic complementarity
    \item[1994] Milgrom \& Shannon: Monotone comparative statics
    \item[2007] Athey: Complementarity under uncertainty
\end{description}
\end{tcolorbox}

\subsection{Consumer Theory: Utility Complementarity}

\subsubsection{The Translation}

\begin{axiom}[Economic Translation E-1]
In consumer theory, the mathematical structure translates as:
\begin{align}
S &\mapsto \text{Set of goods/services } \{1, 2, \ldots, n\} \\
f &\mapsto \text{Utility function } U: \mathbb{R}^n_+ \to \mathbb{R} \\
\gamma_{ij} &\mapsto \frac{\partial^2 U}{\partial x_i \partial x_j} \text{ (Edgeworth complementarity)}
\end{align}
\end{axiom}

\subsubsection{Three Economic Definitions}

\begin{definition}[Edgeworth Complementarity]
Goods $i$ and $j$ are \textbf{Edgeworth complements} if:
\begin{equation}
\gamma^E_{ij} = \frac{\partial^2 U(x)}{\partial x_i \partial x_j} > 0
\end{equation}
Interpretation: Marginal utility of $i$ increases with consumption of $j$.
\end{definition}

\begin{definition}[Hicks-Allen Complementarity]
Goods $i$ and $j$ are \textbf{Hicks-Allen complements} if the compensated cross-price elasticity is negative:
\begin{equation}
\gamma^{HA}_{ij} = -\frac{\partial h_i(p, u)}{\partial p_j} > 0
\end{equation}
where $h_i$ is the Hicksian demand. Interpretation: Compensated demand for $i$ falls when $p_j$ rises.
\end{definition}

\begin{definition}[Gross Complementarity]
Goods $i$ and $j$ are \textbf{gross complements} if:
\begin{equation}
\gamma^G_{ij} = -\frac{\partial x_i(p, m)}{\partial p_j} > 0
\end{equation}
where $x_i$ is Marshallian demand. Interpretation: Demand for $i$ falls when $p_j$ rises.
\end{definition}

\subsubsection{Relationships Between Definitions}

\begin{theorem}[Complementarity Hierarchy]
For twice-differentiable utility with interior solutions:
\begin{enumerate}
    \item Edgeworth $\not\Rightarrow$ Hicks-Allen (income effects matter)
    \item Hicks-Allen $\not\Rightarrow$ Gross (income effects matter)
    \item Under separability: Edgeworth $\Leftrightarrow$ Hicks-Allen
\end{enumerate}
\end{theorem}

\subsubsection{Empirical Estimates: Consumer Complementarity}

\begin{table}[h]
\centering
\caption{Empirical Estimates of Consumer Complementarity}
\label{tab:consumer-gamma}
\renewcommand{\arraystretch}{1.2}
\small
\begin{tabular}{@{}llcl@{}}
\toprule
\textbf{Good Pair} & \textbf{Measure} & \textbf{$\gamma$ Estimate} & \textbf{Source} \\
\midrule
Coffee \& Sugar & Cross-elasticity & $-0.15$ to $-0.30$ & Deaton (1988) \\
Bread \& Butter & Cross-elasticity & $-0.10$ to $-0.25$ & Banks et al. (1997) \\
Cars \& Gasoline & Cross-elasticity & $-0.40$ to $-0.60$ & West (2004) \\
Hardware \& Software & Cross-elasticity & $-0.50$ to $-0.80$ & Brynjolfsson et al. (1996) \\
Leisure \& Entertainment & Edgeworth & $+0.05$ to $+0.20$ & Browning et al. (1985) \\
\bottomrule
\end{tabular}
\end{table}

\subsection{Production Theory: Input Complementarity}

\subsubsection{The Translation}

\begin{axiom}[Production Translation E-2]
In production theory:
\begin{align}
S &\mapsto \text{Set of inputs } \{K, L, M, \ldots\} \\
f &\mapsto \text{Production function } F: \mathbb{R}^n_+ \to \mathbb{R}_+ \\
\gamma_{ij} &\mapsto \frac{\partial^2 F}{\partial x_i \partial x_j} \text{ or } \sigma_{ij} \text{ (elasticity of substitution)}
\end{align}
\end{axiom}

\subsubsection{Functional Forms and Implied Complementarity}

\begin{table}[h]
\centering
\caption{Production Functions and Complementarity}
\label{tab:production-gamma}
\renewcommand{\arraystretch}{1.3}
\begin{tabular}{@{}lll@{}}
\toprule
\textbf{Function} & \textbf{Form} & \textbf{Complementarity $\gamma$} \\
\midrule
Cobb-Douglas & $F = A K^\alpha L^\beta$ & $\gamma_{KL} = \alpha\beta A K^{\alpha-1} L^{\beta-1} > 0$ \\
CES & $F = A[\delta K^\rho + (1-\delta)L^\rho]^{1/\rho}$ & $\sigma = \frac{1}{1-\rho}$; $\rho < 0 \Rightarrow$ complements \\
Leontief & $F = \min\{aK, bL\}$ & Perfect complements ($\sigma = 0$) \\
Linear & $F = aK + bL$ & Perfect substitutes ($\gamma = 0$) \\
Translog & $\ln F = \alpha_0 + \sum \alpha_i \ln x_i + \frac{1}{2}\sum\sum \gamma_{ij} \ln x_i \ln x_j$ & $\gamma_{ij}$ directly estimated \\
\bottomrule
\end{tabular}
\end{table}

\subsubsection{Capital-Skill Complementarity}

\begin{tcolorbox}[colback=green!5!white,colframe=green!75!black,title=\textbf{Classic Finding: Capital-Skill Complementarity}]
\textbf{Griliches (1969):} Physical capital is more complementary with skilled labor than unskilled labor.
\begin{equation}
\gamma_{K,H} > \gamma_{K,L}
\end{equation}
where $H$ = skilled labor, $L$ = unskilled labor.

\textbf{Implications:}
\begin{itemize}[noitemsep]
    \item Capital deepening raises skill premium
    \item Technology adoption favors educated workers
    \item Explains rising wage inequality with technological change
\end{itemize}

\textbf{Empirical estimates:} $\sigma_{KH} \approx 0.5$--$0.7$ vs. $\sigma_{KL} \approx 1.0$--$1.5$
\end{tcolorbox}

\subsection{Game Theory: Strategic Complementarity}

\subsubsection{The Translation}

\begin{axiom}[Game-Theoretic Translation E-3]
In game theory:
\begin{align}
S &\mapsto \text{Strategy profile space } \prod_{i} S_i \\
f &\mapsto \text{Payoff function } \pi_i: \prod_j S_j \to \mathbb{R} \\
\gamma_{ij} &\mapsto \frac{\partial^2 \pi_i}{\partial s_i \partial s_j} \text{ (strategic complementarity)}
\end{align}
\end{axiom}

\begin{definition}[Strategic Complements]
Actions $s_i$ and $s_j$ are \textbf{strategic complements} if:
\begin{equation}
\gamma^{strat}_{ij} = \frac{\partial^2 \pi_i(s_i, s_{-i})}{\partial s_i \partial s_j} > 0
\end{equation}
Interpretation: Player $i$'s marginal return to $s_i$ increases when player $j$ increases $s_j$.
\end{definition}

\subsubsection{Examples of Strategic Complementarity}

\begin{table}[h]
\centering
\caption{Strategic Complementarity Examples}
\label{tab:strategic-gamma}
\renewcommand{\arraystretch}{1.3}
\small
\begin{tabular}{@{}p{3cm}p{3cm}p{2.5cm}p{3.5cm}@{}}
\toprule
\textbf{Context} & \textbf{Strategies} & \textbf{$\gamma$ Sign} & \textbf{Interpretation} \\
\midrule
Price competition (Bertrand) & Prices & $\gamma > 0$ & Higher rival price $\to$ raise own price \\
Quantity competition (Cournot) & Quantities & $\gamma < 0$ & Higher rival quantity $\to$ lower own quantity \\
R\&D investment & R\&D spending & $\gamma > 0$ & Spillovers make R\&D more valuable \\
Network adoption & Adoption decision & $\gamma > 0$ & More adopters $\to$ adoption more valuable \\
Bank runs & Withdrawal & $\gamma > 0$ & Others withdraw $\to$ I should withdraw \\
\bottomrule
\end{tabular}
\end{table}

\subsubsection{Equilibrium Implications}

\begin{theorem}[Milgrom-Roberts 1990]
In games with strategic complements on a lattice:
\begin{enumerate}
    \item The set of Nash equilibria forms a complete lattice
    \item There exist smallest and largest equilibria
    \item Best-response dynamics converge to an equilibrium
    \item Positive shocks shift equilibria upward
\end{enumerate}
\end{theorem}

\subsection{Domain-Specific Axioms: Economics}

\begin{axiom}[E-A1: Diminishing Complementarity]
In most economic contexts, complementarity diminishes at scale:
\begin{equation}
\frac{\partial \gamma_{ij}}{\partial x_i} < 0 \quad \text{and} \quad \frac{\partial \gamma_{ij}}{\partial x_j} < 0
\end{equation}
\end{axiom}

\begin{axiom}[E-A2: Budget Constraint]
Consumer complementarity operates under budget constraint:
\begin{equation}
\sum_i p_i x_i \leq m \quad \Rightarrow \quad \text{trade-offs between complements}
\end{equation}
\end{axiom}

\begin{axiom}[E-A3: Revealed Preference Consistency]
Observed choices must be consistent with complementarity structure:
\begin{equation}
x \succ y \text{ and } \gamma_{xy} > 0 \Rightarrow \text{bundle containing both } x,y \text{ chosen over either alone}
\end{equation}
\end{axiom}

% ============================================================================
% PART III: PSYCHOLOGICAL FOUNDATIONS
% ============================================================================
\section{Psychological Foundations of Complementarity}
\label{sec:psychology}

\subsection{Behavioral Complementarity}

\subsubsection{The Translation}

\begin{axiom}[Psychological Translation P-1]
In behavioral psychology:
\begin{align}
S &\mapsto \text{Set of behaviors } \{B_1, B_2, \ldots, B_n\} \\
f &\mapsto \text{Outcome function } O: 2^S \to \mathbb{R} \text{ (well-being, performance)} \\
\gamma_{ij} &\mapsto \text{Behavioral interaction effect}
\end{align}
\end{axiom}

\subsubsection{Types of Behavioral Complementarity}

\begin{definition}[Behavioral Synergy]
Behaviors $B_i$ and $B_j$ exhibit \textbf{behavioral synergy} if:
\begin{equation}
\gamma^{beh}_{ij} = O(B_i \cap B_j) - O(B_i) - O(B_j) + O(\emptyset) > 0
\end{equation}
Example: Exercise + healthy diet $>$ exercise alone + diet alone.
\end{definition}

\begin{definition}[Habit Complementarity]
Habits $H_i$ and $H_j$ are \textbf{habit complements} if establishing one makes the other easier:
\begin{equation}
P(H_j | H_i) > P(H_j | \neg H_i)
\end{equation}
Example: Morning routine (wake early + exercise + healthy breakfast).
\end{definition}

\subsubsection{Empirical Evidence: Behavior Clusters}

\begin{table}[h]
\centering
\caption{Empirical Behavioral Complementarities}
\label{tab:behavior-gamma}
\renewcommand{\arraystretch}{1.2}
\small
\begin{tabular}{@{}p{4cm}p{2.5cm}p{2cm}l@{}}
\toprule
\textbf{Behavior Pair} & \textbf{Outcome} & \textbf{$\gamma$ Effect} & \textbf{Source} \\
\midrule
Exercise + Diet & Weight loss & $+15\%$ synergy & Franz et al. (2007) \\
Smoking cessation + Exercise & Quit success & $+25\%$ synergy & Marcus et al. (1999) \\
Mindfulness + CBT & Depression & $+20\%$ synergy & Teasdale et al. (2000) \\
Sleep + Exercise & Cognitive function & $+30\%$ synergy & Scullin \& Bliwise (2015) \\
Social support + Self-efficacy & Goal achievement & $+40\%$ synergy & Bandura (1997) \\
\bottomrule
\end{tabular}
\end{table}

\subsection{Cognitive Complementarity}

\subsubsection{Dual-Process Complementarity}

\begin{tcolorbox}[colback=purple!5!white,colframe=purple!75!black,title=\textbf{Kahneman's System 1/System 2 Complementarity}]
Cognitive processing involves complementary systems:
\begin{equation}
\text{Performance} = f(S_1, S_2) \quad \text{with} \quad \gamma_{S_1,S_2} \text{ context-dependent}
\end{equation}

\textbf{Complementarity conditions:}
\begin{itemize}[noitemsep]
    \item \textbf{High $\gamma$:} Complex decisions requiring both intuition and analysis
    \item \textbf{Low $\gamma$:} Simple routine decisions (System 1 sufficient)
    \item \textbf{Negative $\gamma$:} Overthinking simple problems (System 2 interference)
\end{itemize}
\end{tcolorbox}

\subsubsection{Skill Complementarity}

\begin{definition}[Cognitive Skill Complementarity]
Skills $K_i$ and $K_j$ are \textbf{cognitive complements} if:
\begin{equation}
\gamma^{cog}_{ij} = \frac{\partial^2 \text{Performance}}{\partial K_i \partial K_j} > 0
\end{equation}
\end{definition}

\begin{table}[h]
\centering
\caption{Cognitive Skill Complementarities}
\label{tab:cognitive-gamma}
\renewcommand{\arraystretch}{1.2}
\begin{tabular}{@{}llcl@{}}
\toprule
\textbf{Skill Pair} & \textbf{Domain} & \textbf{$\gamma$ Sign} & \textbf{Evidence} \\
\midrule
Verbal + Quantitative & Academic & $\gamma > 0$ & Deary et al. (2007) \\
Creativity + Technical & Innovation & $\gamma > 0$ & Amabile (1996) \\
Emotional + Analytical & Leadership & $\gamma > 0$ & Goleman (1995) \\
Memory + Processing speed & Problem-solving & $\gamma > 0$ & Salthouse (2004) \\
\bottomrule
\end{tabular}
\end{table}

\subsection{Motivational Complementarity}

\subsubsection{Intrinsic-Extrinsic Interaction}

\begin{theorem}[Crowding Effects]
Intrinsic motivation ($M_I$) and extrinsic incentives ($M_E$) interact:
\begin{equation}
\gamma_{IE} = \frac{\partial^2 \text{Effort}}{\partial M_I \partial M_E}
\end{equation}
\begin{itemize}
    \item \textbf{Crowding-in} ($\gamma > 0$): Extrinsic rewards enhance intrinsic motivation
    \item \textbf{Crowding-out} ($\gamma < 0$): Extrinsic rewards undermine intrinsic motivation
\end{itemize}
\end{theorem}

\textbf{Empirical findings (Deci et al., 1999 meta-analysis):}
\begin{itemize}[noitemsep]
    \item Performance-contingent rewards: $\gamma \approx -0.15$ (crowding-out)
    \item Engagement-contingent rewards: $\gamma \approx +0.05$ (slight crowding-in)
    \item Unexpected rewards: $\gamma \approx 0$ (neutral)
\end{itemize}

\subsection{Domain-Specific Axioms: Psychology}

\begin{axiom}[P-A1: Bounded Rationality]
Psychological complementarity operates under cognitive constraints:
\begin{equation}
\text{Effective } \gamma^{eff}_{ij} = \gamma_{ij} \cdot A(\cdot)
\end{equation}
where $A(\cdot) \in [0,1]$ is the awareness function (see CORE-AWARE).
\end{axiom}

\begin{axiom}[P-A2: Temporal Decay]
Behavioral complementarity decays without reinforcement:
\begin{equation}
\gamma_{ij}(t) = \gamma_{ij}(0) \cdot e^{-\lambda t} \quad \text{unless reinforced}
\end{equation}
\end{axiom}

\begin{axiom}[P-A3: Context Sensitivity]
Psychological $\gamma$ is highly context-dependent:
\begin{equation}
\gamma_{ij}(\Psi_1) \neq \gamma_{ij}(\Psi_2) \quad \text{for different contexts } \Psi
\end{equation}
\end{axiom}

% ============================================================================
% PART IV: MANAGEMENT FOUNDATIONS
% ============================================================================
\section{Management Foundations of Complementarity}
\label{sec:management}

\subsection{Organizational Complementarity}

\subsubsection{The Translation}

\begin{axiom}[Management Translation M-1]
In organizational theory:
\begin{align}
S &\mapsto \text{Set of practices/resources } \{P_1, P_2, \ldots, P_n\} \\
f &\mapsto \text{Performance function } \Pi: 2^S \to \mathbb{R} \\
\gamma_{ij} &\mapsto \text{Organizational fit / strategic alignment}
\end{align}
\end{axiom}

\subsubsection{The Milgrom-Roberts Framework}

\begin{tcolorbox}[colback=blue!5!white,colframe=blue!75!black,title=\textbf{Milgrom \& Roberts (1990, 1995): Organizational Complementarities}]
Modern manufacturing involves \textbf{clusters of complementary practices}:
\begin{equation}
\Pi(\mathbf{x}) = \sum_i \alpha_i x_i + \sum_{i < j} \gamma_{ij} x_i x_j + \sum_{i < j < k} \gamma_{ijk} x_i x_j x_k + \ldots
\end{equation}

\textbf{Key complementary clusters:}
\begin{itemize}[noitemsep]
    \item Flexible manufacturing + worker training + supplier relations
    \item IT investment + organizational restructuring + skill upgrading
    \item Quality focus + employee empowerment + customer orientation
\end{itemize}

\textbf{Implication:} Piecemeal adoption fails; complementary practices must be adopted together.
\end{tcolorbox}

\subsection{Strategic Fit}

\begin{definition}[Strategic Complementarity]
Strategies $S_i$ and $S_j$ exhibit \textbf{strategic fit} if:
\begin{equation}
\gamma^{strat}_{ij} = \Pi(S_i, S_j) - \Pi(S_i, \neg S_j) - \Pi(\neg S_i, S_j) + \Pi(\neg S_i, \neg S_j) > 0
\end{equation}
\end{definition}

\subsubsection{Porter's Activity Systems}

\begin{table}[h]
\centering
\caption{Porter's Activity System Complementarities}
\label{tab:porter-gamma}
\renewcommand{\arraystretch}{1.2}
\small
\begin{tabular}{@{}p{2.5cm}p{4cm}p{2cm}p{3cm}@{}}
\toprule
\textbf{Company} & \textbf{Complementary Activities} & \textbf{$\gamma$ Type} & \textbf{Competitive Effect} \\
\midrule
Southwest Airlines & Point-to-point + No meals + Fast turnaround & High $\gamma$ & Cost leadership \\
IKEA & Self-service + Flat-pack + Suburban stores & High $\gamma$ & Cost + convenience \\
Vanguard & Index funds + Low fees + Direct sales & High $\gamma$ & Cost leadership \\
Apple & Design + Ecosystem + Retail stores & High $\gamma$ & Differentiation \\
\bottomrule
\end{tabular}
\end{table}

\subsection{Human Resource Complementarities}

\begin{theorem}[Ichniowski et al., 1997]
\textbf{HR practice clusters} exhibit strong complementarity:
\begin{equation}
\Pi(\text{All practices}) >> \sum_i \Pi(\text{Practice}_i)
\end{equation}

\textbf{Complementary HR practices:}
\begin{enumerate}[noitemsep]
    \item Incentive pay + Employment security
    \item Training + Participation
    \item Flexible job design + Information sharing
\end{enumerate}

\textbf{Empirical finding:} Steel mills with full system: +7\% productivity vs. traditional.
\end{theorem}

\subsection{IT-Organization Complementarity}

\begin{tcolorbox}[colback=green!5!white,colframe=green!75!black,title=\textbf{Brynjolfsson \& Hitt (2000): IT and Organizational Capital}]
IT investment ($IT$) is complementary with organizational practices ($ORG$):
\begin{equation}
\gamma_{IT,ORG} = \frac{\partial^2 \Pi}{\partial IT \partial ORG} > 0
\end{equation}

\textbf{Components of organizational capital:}
\begin{itemize}[noitemsep]
    \item Decentralized decision-making
    \item Skilled workforce
    \item Process redesign
    \item Supplier/customer integration
\end{itemize}

\textbf{Empirical finding:} IT + organizational change: 5-year return $>$ 80\%; IT alone: return $<$ 20\%.
\end{tcolorbox}

\subsection{Domain-Specific Axioms: Management}

\begin{axiom}[M-A1: Implementation Lag]
Organizational complementarity requires time to materialize:
\begin{equation}
\gamma^{eff}_{ij}(t) = \gamma^{pot}_{ij} \cdot (1 - e^{-\lambda t})
\end{equation}
where $\lambda$ is the implementation rate.
\end{axiom}

\begin{axiom}[M-A2: Switching Costs]
Complementary systems create lock-in:
\begin{equation}
\text{Switching cost} = \sum_{i,j} \gamma_{ij} \cdot \mathbf{1}[\text{both } i,j \text{ adopted}]
\end{equation}
\end{axiom}

\begin{axiom}[M-A3: Adoption Indivisibility]
Complementary practices are optimally adopted together:
\begin{equation}
\gamma_{ij} > c_i + c_j \Rightarrow \text{Adopt both or neither}
\end{equation}
where $c_i$ is the adoption cost of practice $i$.
\end{axiom}

% ============================================================================
% PART V: CROSS-DOMAIN COMPARISONS
% ============================================================================
\section{Cross-Domain Comparison}
\label{sec:comparison}

\subsection{Structural Similarities}

\begin{table}[h]
\centering
\caption{Cross-Domain Complementarity Comparison}
\label{tab:cross-domain}
\renewcommand{\arraystretch}{1.3}
\small
\begin{tabular}{@{}p{2.5cm}p{3.5cm}p{3.5cm}p{3.5cm}@{}}
\toprule
\textbf{Aspect} & \textbf{Economics} & \textbf{Psychology} & \textbf{Management} \\
\midrule
\textbf{Elements} & Goods, inputs & Behaviors, cognitions & Practices, resources \\
\textbf{Value function} & Utility, production & Well-being, performance & Firm value \\
\textbf{Measurement} & Prices, quantities & Surveys, experiments & Productivity, profits \\
\textbf{Typical $\gamma$ range} & $-0.5$ to $+0.5$ & $-0.3$ to $+0.4$ & $0$ to $+0.8$ \\
\textbf{Time scale} & Instantaneous to years & Days to months & Months to years \\
\textbf{Key constraint} & Budget & Cognitive capacity & Resources, time \\
\bottomrule
\end{tabular}
\end{table}

\subsection{Domain-Specific Measurement Challenges}

\begin{table}[h]
\centering
\caption{Measurement Approaches by Domain}
\label{tab:measurement}
\renewcommand{\arraystretch}{1.3}
\small
\begin{tabular}{@{}p{2.5cm}p{4.5cm}p{5cm}@{}}
\toprule
\textbf{Domain} & \textbf{Preferred Methods} & \textbf{Key Challenges} \\
\midrule
\textbf{Economics} & Demand estimation, production function estimation, experiments & Identification, endogeneity, functional form \\
\textbf{Psychology} & RCTs, factorial designs, longitudinal studies & Sample size, external validity, measurement error \\
\textbf{Management} & Case studies, panel data, natural experiments & Selection bias, unobserved heterogeneity, causality \\
\bottomrule
\end{tabular}
\end{table}

\subsection{Universal Properties Across Domains}

\begin{theorem}[Cross-Domain Universals]
The following properties hold across all three domains:
\begin{enumerate}
    \item \textbf{Superadditivity when $\gamma > 0$:} Combined effect exceeds sum of parts
    \item \textbf{Adoption clustering:} Complementary elements tend to co-occur
    \item \textbf{Amplification:} Shocks to one element affect complementary elements
    \item \textbf{Path dependence:} Order of adoption matters when $\gamma \neq 0$
\end{enumerate}
\end{theorem}

% ============================================================================
% PART VI: PARAMETER ESTIMATION ACROSS DOMAINS
% ============================================================================
\section{Parameter Estimation Across Domains}
\label{sec:parameters}

\subsection{Estimation Hierarchy}

Following CORE-WHERE (Appendix BBB), we establish domain-specific estimation priorities:

\begin{table}[h]
\centering
\caption{Parameter Estimation by Domain and Method}
\label{tab:estimation}
\renewcommand{\arraystretch}{1.3}
\small
\begin{tabular}{@{}lccc@{}}
\toprule
\textbf{Method} & \textbf{Economics} & \textbf{Psychology} & \textbf{Management} \\
\midrule
Meta-analysis & Available & Rich & Limited \\
Natural experiments & Good & Limited & Good \\
RCTs & Some & Many & Few \\
Structural estimation & Strong tradition & Emerging & Limited \\
Survey/self-report & Secondary & Primary & Mixed \\
LLM Monte Carlo & Supplementary & Useful & Valuable \\
\bottomrule
\end{tabular}
\end{table}

\subsection{Recommended $\gamma$ Priors by Domain}

\begin{table}[h]
\centering
\caption{Recommended Prior Distributions for $\gamma$}
\label{tab:priors}
\renewcommand{\arraystretch}{1.3}
\begin{tabular}{@{}llll@{}}
\toprule
\textbf{Domain} & \textbf{Context} & \textbf{Prior} & \textbf{Justification} \\
\midrule
\multirow{2}{*}{Economics} & Consumer goods & $\gamma \sim N(0, 0.2^2)$ & Most pairs independent \\
 & Production inputs & $\gamma \sim N(0.1, 0.15^2)$ & Capital-labor complement \\
\multirow{2}{*}{Psychology} & Health behaviors & $\gamma \sim N(0.15, 0.1^2)$ & Behavior clustering \\
 & Cognitive skills & $\gamma \sim N(0.1, 0.1^2)$ & Positive manifold \\
\multirow{2}{*}{Management} & HR practices & $\gamma \sim N(0.2, 0.15^2)$ & Strong clustering \\
 & IT + Organization & $\gamma \sim N(0.3, 0.2^2)$ & High complementarity \\
\bottomrule
\end{tabular}
\end{table}

% ============================================================================
% PART VII: BRIDGE TO EBF (LAYER 4)
% ============================================================================
\section{Bridge to EBF Application (Layer 4)}
\label{sec:bridge}

\subsection{The EBF Translation}

\begin{tcolorbox}[colback=red!5!white,colframe=red!75!black,title=\textbf{From Domain Foundations to EBF}]
Layer 4 (Appendix XIII) will apply these domain foundations to behavioral change:
\begin{equation}
\gamma^{EBF}_{ij} = \text{Complementarity between interventions } i \text{ and } j
\end{equation}

\textbf{Key EBF elements:}
\begin{itemize}[noitemsep]
    \item Interventions as elements (nudges, incentives, information, defaults)
    \item Behavioral outcomes as value function
    \item FEPSDE dimensions interact via $\gamma$
    \item Context $\Psi$ modulates all $\gamma$ values
\end{itemize}
\end{tcolorbox}

\subsection{Required Inputs for Layer 4}

\begin{enumerate}
    \item \textbf{From Economics:} Consumer choice models, strategic interaction framework
    \item \textbf{From Psychology:} Behavioral synergy estimates, habit complementarity
    \item \textbf{From Management:} Implementation dynamics, adoption clustering
\end{enumerate}

% ============================================================================
% GLOSSARY
% ============================================================================
\section{Glossary}
\label{sec:glossary-domain}

For complete glossary, see Appendix G (REF-GLOSSARY).

\begin{description}[noitemsep]
    \item[Domain Translation] Mapping from mathematical structure to domain-specific interpretation
    \item[Edgeworth Complementarity] $\partial^2 U / \partial x_i \partial x_j > 0$
    \item[Hicks-Allen Complementarity] Negative compensated cross-price effect
    \item[Strategic Complementarity] $\partial^2 \pi_i / \partial s_i \partial s_j > 0$
    \item[Behavioral Synergy] Combined behavioral effect exceeds sum of individual effects
    \item[Organizational Fit] Complementarity among management practices
    \item[Crowding Effect] Interaction between intrinsic and extrinsic motivation
\end{description}

% ============================================================================
% REFERENCES
% ============================================================================
\section{References}
\label{sec:references-domain}

\subsection{Economics (20 papers)}

\begin{enumerate}
    \item Edgeworth, F.Y. (1881). \textit{Mathematical Psychics}. Kegan Paul.
    \item Hicks, J.R. \& Allen, R.G.D. (1934). A reconsideration of the theory of value. \textit{Economica}, 1(1), 52--76.
    \item Samuelson, P.A. (1974). Complementarity: An essay on the 40th anniversary of the Hicks-Allen revolution. \textit{Journal of Economic Literature}, 12(4), 1255--1289.
    \item Topkis, D.M. (1978). Minimizing a submodular function on a lattice. \textit{Operations Research}, 26(2), 305--321.
    \item Milgrom, P. \& Roberts, J. (1990). Rationalizability, learning, and equilibrium. \textit{Econometrica}, 58(6), 1255--1277.
    \item Milgrom, P. \& Shannon, C. (1994). Monotone comparative statics. \textit{Econometrica}, 62(1), 157--180.
    \item Vives, X. (1990). Nash equilibrium with strategic complementarities. \textit{Journal of Mathematical Economics}, 19(3), 305--321.
    \item Athey, S. (2002). Monotone comparative statics under uncertainty. \textit{QJE}, 117(1), 187--223.
    \item Griliches, Z. (1969). Capital-skill complementarity. \textit{Review of Economics and Statistics}, 51(4), 465--468.
    \item Krusell, P. et al. (2000). Capital-skill complementarity and inequality. \textit{Econometrica}, 68(5), 1029--1053.
    \item Deaton, A. \& Muellbauer, J. (1980). \textit{Economics and Consumer Behavior}. Cambridge.
    \item Banks, J., Blundell, R. \& Lewbel, A. (1997). Quadratic Engel curves and consumer demand. \textit{Review of Economics and Statistics}, 79(4), 527--539.
    \item Brynjolfsson, E. \& Kemerer, C.F. (1996). Network externalities in microcomputer software. \textit{Management Science}, 42(12), 1627--1647.
    \item West, S.E. (2004). Distributional effects of alternative vehicle pollution control policies. \textit{Journal of Public Economics}, 88(3-4), 735--757.
    \item Browning, M. (1991). A simple nonadditive preference structure for models of household behavior. \textit{Journal of Political Economy}, 99(3), 607--637.
    \item Quah, J.K.-H. (2007). The comparative statics of constrained optimization. \textit{Econometrica}, 75(2), 401--431.
    \item Echenique, F. (2002). Comparative statics by adaptive dynamics. \textit{Econometrica}, 70(2), 833--844.
    \item Cooper, R. \& John, A. (1988). Coordinating coordination failures in Keynesian models. \textit{QJE}, 103(3), 441--463.
    \item Bulow, J.I., Geanakoplos, J.D. \& Klemperer, P.D. (1985). Multimarket oligopoly: Strategic substitutes and complements. \textit{Journal of Political Economy}, 93(3), 488--511.
    \item Diamond, P.A. (1982). Aggregate demand management in search equilibrium. \textit{Journal of Political Economy}, 90(5), 881--894.
\end{enumerate}

\subsection{Psychology (15 papers)}

\begin{enumerate}
    \setcounter{enumi}{20}
    \item Kahneman, D. (2011). \textit{Thinking, Fast and Slow}. Farrar, Straus and Giroux.
    \item Bandura, A. (1997). \textit{Self-Efficacy: The Exercise of Control}. Freeman.
    \item Deci, E.L., Koestner, R. \& Ryan, R.M. (1999). A meta-analytic review of experiments examining the effects of extrinsic rewards on intrinsic motivation. \textit{Psychological Bulletin}, 125(6), 627--668.
    \item Franz, M.J. et al. (2007). Weight-loss outcomes: A systematic review and meta-analysis. \textit{Journal of the American Dietetic Association}, 107(10), 1755--1767.
    \item Marcus, B.H. et al. (1999). The efficacy of exercise as an aid for smoking cessation in women. \textit{Archives of Internal Medicine}, 159(11), 1229--1234.
    \item Teasdale, J.D. et al. (2000). Prevention of relapse/recurrence in major depression by mindfulness-based cognitive therapy. \textit{Journal of Consulting and Clinical Psychology}, 68(4), 615--623.
    \item Scullin, M.K. \& Bliwise, D.L. (2015). Sleep, cognition, and normal aging. \textit{Perspectives on Psychological Science}, 10(1), 97--137.
    \item Deary, I.J. et al. (2007). Intelligence and educational achievement. \textit{Intelligence}, 35(1), 13--21.
    \item Amabile, T.M. (1996). \textit{Creativity in Context}. Westview Press.
    \item Goleman, D. (1995). \textit{Emotional Intelligence}. Bantam Books.
    \item Salthouse, T.A. (2004). What and when of cognitive aging. \textit{Current Directions in Psychological Science}, 13(4), 140--144.
    \item Prochaska, J.O. \& DiClemente, C.C. (1983). Stages and processes of self-change of smoking. \textit{Journal of Consulting and Clinical Psychology}, 51(3), 390--395.
    \item Michie, S., van Stralen, M.M. \& West, R. (2011). The behaviour change wheel. \textit{Implementation Science}, 6(1), 42.
    \item Wood, W. \& Rünger, D. (2016). Psychology of habit. \textit{Annual Review of Psychology}, 67, 289--314.
    \item Sheeran, P. \& Webb, T.L. (2016). The intention-behavior gap. \textit{Social and Personality Psychology Compass}, 10(9), 503--518.
\end{enumerate}

\subsection{Management (15 papers)}

\begin{enumerate}
    \setcounter{enumi}{35}
    \item Milgrom, P. \& Roberts, J. (1990). The economics of modern manufacturing. \textit{American Economic Review}, 80(3), 511--528.
    \item Milgrom, P. \& Roberts, J. (1995). Complementarities and fit: Strategy, structure, and organizational change. \textit{Journal of Accounting and Economics}, 19(2-3), 179--208.
    \item Ichniowski, C., Shaw, K. \& Prennushi, G. (1997). The effects of human resource management practices on productivity. \textit{American Economic Review}, 87(3), 291--313.
    \item Brynjolfsson, E. \& Hitt, L.M. (2000). Beyond computation: Information technology, organizational transformation and business performance. \textit{Journal of Economic Perspectives}, 14(4), 23--48.
    \item Porter, M.E. (1996). What is strategy? \textit{Harvard Business Review}, 74(6), 61--78.
    \item Siggelkow, N. (2002). Evolution toward fit. \textit{Administrative Science Quarterly}, 47(1), 125--159.
    \item Rivkin, J.W. (2000). Imitation of complex strategies. \textit{Management Science}, 46(6), 824--844.
    \item Bloom, N. \& Van Reenen, J. (2007). Measuring and explaining management practices across firms and countries. \textit{Quarterly Journal of Economics}, 122(4), 1351--1408.
    \item Huselid, M.A. (1995). The impact of human resource management practices on turnover, productivity, and corporate financial performance. \textit{Academy of Management Journal}, 38(3), 635--672.
    \item MacDuffie, J.P. (1995). Human resource bundles and manufacturing performance. \textit{Industrial and Labor Relations Review}, 48(2), 197--221.
    \item Teece, D.J., Pisano, G. \& Shuen, A. (1997). Dynamic capabilities and strategic management. \textit{Strategic Management Journal}, 18(7), 509--533.
    \item Barney, J. (1991). Firm resources and sustained competitive advantage. \textit{Journal of Management}, 17(1), 99--120.
    \item Ennen, E. \& Richter, A. (2010). The whole is more than the sum of its parts—or is it? A review of the empirical literature on complementarities in organizations. \textit{Journal of Management}, 36(1), 207--233.
    \item Cassiman, B. \& Veugelers, R. (2006). In search of complementarity in innovation strategy. \textit{Management Science}, 52(1), 68--82.
    \item Arora, A. \& Gambardella, A. (1990). Complementarity and external linkages. \textit{Journal of Industrial Economics}, 38(4), 361--379.
\end{enumerate}

% ============================================================================
% SUMMARY
% ============================================================================
\section{Summary}
\label{sec:summary}

This appendix established the domain foundations of complementarity, bridging:
\begin{itemize}[noitemsep]
    \item \textbf{Layer 1} (Appendix X): Pure mathematical theory
    \item \textbf{Layer 2} (Appendix XI): Canonical mathematical models
    \item \textbf{Layer 3} (This Appendix): Domain-specific interpretations
    \item \textbf{Layer 4} (Appendix XIII): EBF behavioral application
\end{itemize}

\textbf{Key contributions:}
\begin{enumerate}
    \item Established translation maps for Economics, Psychology, Management
    \item Documented domain-specific axioms and constraints
    \item Compiled empirical $\gamma$ estimates from each domain
    \item Identified measurement challenges and recommended priors
    \item Created bridge to EBF application layer
\end{enumerate}

\end{chapter}
